\newpage

\section{我的姐姐李新瑩}
\label{sister}

在社會上行走久了，感覺很多人不像人，更像是畜生：有人是雞、有人是鴨、有人是牛馬。有的畜生向上修煉成精，比如我的父親，地產商李松堅，我媽媽很小的時候就告訴我他其實是豬精\footnote{我媽媽對李松堅的稱呼是潮汕話「豬哥精」。}。而有的畜生則向下繼續墮落，比如我姐姐李新瑩，給雞當狗，成為畜生的畜生。

\ 

我姐姐的畜生道，還要從1992年的那個夏夜開始說起。

1989年，我的姐姐出生在廣東汕頭。本來她是拿着宗馥莉般的大女主劇本的：家族長女，爺爺是鎮長，外公是校長，奶奶和外婆是泰國和臺灣歸國家庭的千金，媽媽是體制內中學教師，而我的父親，在90年代初就是資產過億的大老闆了。可是，沒過幾年，父親就出軌了。他出軌的小三，是一個在廣東打工的二婚帶男娃的狐狸精，名字叫做凌菲菲。

很快，父親就和狐狸精同居了。為了討狐狸精開心，他把我姐姐接過去和狐狸精同住，跟狐狸精一起睡，還讓我姐姐叫狐狸精媽媽。狐狸精有肺結核。很快，我的姐姐也被傳染了肺結核。當時我姐姐才3歲左右，身體抵抗力很差。很快她就肺結核發病，高燒，昏迷。當時醫生把我姐姐抱起來，她的頭、手都是直接耷拉下去，一點意識都沒有。醫生說，這個小女孩可能救不活了。我的父親看到我姐姐這個樣子，突然良心發現。他抱着我的姐姐，用盡所有方法搶救。用了超過成人劑量幾倍的抗生素和激素，終於把我姐姐搶救回來了。

在那個1992年的夏夜，三歲的姐姐從瀕死狀態蘇醒過來，迷迷蒙蒙之中，睜開雙眼，是父親偉岸的身影，和充滿父愛的呼喚。她又閉上了雙眼，安靜地睡去。但是那一夜那炙熱的父愛，已經永遠地植入了她心中的最深處。

但是我的姐姐不知道的是，父親對她的父愛，僅限於她差點被狐狸精害死的那一夜。很快，父親就嫌她煩人，把她送回給了我媽媽，和狐狸精雙宿雙棲去了。送回去的時候，父親並沒有把我姐姐得肺結核瀕死的事情告知我媽媽，我和我媽媽和我外公都感染了結核菌。我和媽媽抵抗力比較強，沒有發展成肺結核。但是我外公肺結核發病，留下了嚴重的後遺症，最後也因為肺部問題去世了。至於我姐姐，在被送回來之後，身體一直很差。康復後半邊臉都是黑點，輕微毀容。而且不知道是因為高燒昏厥、抗生素還是激素，我姐姐受到了永久性的智力損傷，成為家族最愚笨，成績最差的一個。\footnote{30年來，李松堅和凌菲菲一直隱瞞我姐姐肺炎瀕死的事，而且也隱瞞故意把肺結核傳染給我的事情，並逃逸定居上海至今。我的姐姐和您一樣，也是從我這裏第一次知道這件事。}

\ 

豬精和狐狸精雙宿雙飛之後，轉移了所有資產，並僱傭黑社會威脅欺負我們：監聽電話、扔死老鼠、辱罵、威脅。我媽媽最後害怕黑社會，也怕我姐姐撫養權被搶走被搞死，被迫簽下離婚協議，只拿到千分之一左右的資產。

可是，偶蹄目的豬，怎麼敵得過食肉目的狐狸呢？豬精因為陳良宇案挪用社保基金，被紀委調查的時候。狐狸精毅然決然地背刺，在豬精消失的時候，背刺起訴，把豬精靠黑社會從我媽媽身上搶奪來的資產，分走一半。包括上海明園集團的一半股份和上海的多套別墅房產。而他最愛的繼子，也不認他了。

豬精被背刺之後，一面假惺惺把我和姐姐接到上海，一面深入反思：以後再也不能找食肉動物了！豬腦子轉了一圈又一圈，還是叫雞好——雞，不僅不是食肉動物，連哺乳動物都不是！於是，他到上海音樂學院，認識了一個龜公，叫做徐沛東。徐沛東當時是上海音樂學院的教授。他在他的藝術生裏面，找了一個願意做雞的，叫做許潔。龜公哼哧哼哧把雞背在身上，獻給了豬精。為了豬精開心，龜公還利用他教授和校董的身份和豬精的錢，賄賂學校，給許潔買了一個教授的身份，變成了教授雞。後來，教授雞在豬精接近60歲高齡的時候，給豬精生了個雄性雞仔，叫做李新權。

有了弟弟之後，我和姐姐的命運就被安排好了。因為我不僅美貌還聰明，獲得了中國電腦奧賽第二名和美國電腦奧賽第一名，我爸就逼我放棄麻省理工的錄取轉去清華，希望我能和領導的子女和親。至於被狐狸精燒壞腦袋的姐姐，出嫁也收益不大，於是我父親把她留在身邊當下人。

2021年我回國，豬精把我和姐姐叫到公司吃飯。他為了給自己小時候拋妻棄子和謀殺子女未遂辯護，對我們說：「你們年輕人，不應該談戀愛，不應該結婚，應該奮鬥事業。等到你們50歲左右才結婚最好。」當時我震驚得不行：50歲我和姐姐都絕經了，還結個屁婚。我驚恐地望向姐姐，結果姐姐專心致志，聽得津津有味。當時我姐姐做一個裝修公司，效益不好關門了。而豬精也被騙，投資了一個所謂高科技塑膠廠，血本無歸。於是豬精努力轉了一下他的豬腦袋，用發現新大陸的語氣，對我們說教道：「高科技，都是騙人的！還是傳統行業好。」於是，他叫我姐姐去開店賣潮汕牛肉火鍋。至於我，作為清華姚班畢業的高材生，應該學以致用，在淘寶開店賣牛肉丸。

雖然一直以來，我一直把豬精的話當豬叫，但是這段發言還是讓我下巴都驚掉了。但是更讓我震驚的是，我姐還真的聽了！在37虛歲的年紀，孑然一身，在潮汕牛肉火鍋店哼哧哼哧證明自己。每天給教授雞當狗，陪她逛街買包。還到處罵我不肯賣牛肉丸和陪教授雞逛街，實在是太不孝順了。

半年前，豬精房地產項目爛尾了，資金鏈斷裂。他聽說我自己賺錢了，於是跑到深圳我公司拜訪我，想跟我借錢。他對我說，他還是喜歡兒子多一點，但是因為我賺到錢，也喜歡我。他說了一堆好聽話，然後送了我一塊2萬的手錶，再把我的100多萬的邁巴赫開走了。我姐姐聽說後，更痛苦了，她不明白，為什麼她聽話賣牛肉火鍋，給雞當狗，做畜生的畜生，為什麼爸爸還是愛我多過愛她。\footnote{我姐姐不知道李松堅借了我的邁巴赫拒絕歸還的事情，希望她知道我被李松堅爆金幣之後，能不那麼嫉妒我，開心一點。}

\ 

每次想到我從小相依為命的姐姐一輩子給NPD豬精父親當血包的人生，我都忍不住流淚。我有時候會想，豬精如果不是這種10年以上，死緩以下的道德觀，也許我姐姐就不會這麼痛苦吧。如果父親不出軌，不犯故意傳播甲類傳染病罪、故意殺人未遂罪、濫用職權罪、組織黑社會性質組織罪等罪，我和姐姐現在應該正自豪地開着父親80年代買的豐田皇冠皇家沙龍，發抖音炫耀。而如果父親能突破犯罪上限，和凌菲菲一起把我姐姐肺結核搞死，一了百了，她也就不用忍受後來的痛苦。或者他能夠突破父女的倫理禁忌，也許他就會像愛狐狸精和雞一樣愛我姐姐，而我姐姐也不會一輩子缺愛了。

我從小就沒有爹，沒有感受過父愛，也因此也自由自在習慣了，實在不能理解我的姐姐為什麼要給雞當狗，做畜生的畜生。或許是那個夏夜的父愛是那麼刻骨銘心，值得讓她終其一生苦苦追尋。

——亦或許，在那個夏夜，她腦子燒壞了吧。

\ \\
李新野\\
2025年7月19日

\ 

--------------------------------------------------------------------------

\

\textit{李新野網誌原文連結：https://sinyalee.com/blog/?p=1018}

\newpage

\section{我的大姨李麗麗}
\label{aunt}
我的大姨媽李麗麗，是汕頭金山中學之前的校長。她也是我遇到過的最矛盾、最扭曲的人。

\ \\
李麗麗是我這輩子見過的最清廉、最偉大的幹部。她最引以自豪的，是她的作為金山中學校長的副處級編制。她最崇拜的，是級別比她高的人。當年她作為人大代表的時候，和主席握了手，她興奮到連續幾天不洗手。過了幾年，還是在金山中學的全校大會上說當年和主席握手的點點滴滴，還說她不洗手的事情。

然而李麗麗最仇恨的對象，第一名是我媽媽，第二名是我爸爸。在她眼中，我媽是一個連黨員都不是的大專生、群眾，是比她低很多等的。結果卻嫁給了我爸爸有錢人，生活比她好很多倍。而我爸爸，是一個連大學都沒讀過的半文盲，結果卻因為很早就做貿易和開工廠，賺到了她1000輩子都賺不到的錢。一個她眼中的底層，哦不，一個她眼中應該下地獄的生意人+家庭主婦家庭，生活卻過得比她好，是她最最最嫉妒和仇恨的。

\ \\
哦，忘了說，她最仇恨的對象，也是她最大的收入來源。

2007年，中國石油上市A股，李麗麗家全副身家幾十萬重倉中國石油股票。很快血本無歸。結果沒過多久，李麗麗就搬新家，裝修新房子了。猜猜為什麼李麗麗能那麼快翻身？因為有我媽媽源源不斷的給她錢。

所以，李麗麗平時對着我的父母，還是把她內心深處的仇恨掩蓋起來，充滿尊重。但是，大人之間的面具，在小孩身上就撕扯得稀爛了。

在我初二的時候，我和我表姐，也是李麗麗的女兒，還有其他親戚一起在雲南玩了一個多星期。我表姐每天指着我鼻子罵，排擠我，搶我的東西，罵我是邪惡的商人、罪犯的兒子。當時我被欺負到崩潰大哭。一個未成年的女孩有這種思想，肯定是李麗麗和她老公日復一日灌輸的結果。（這也是我初三突然努力讀書的原因。就是為了考到廣州不要和李麗麗一起。）

後來我研究生畢業再次見到我表姐的時候。當時她在上海我爸爸安排的工作中繼續靠我父母接濟。而我在美國靠自己工作，年收入一百萬人民幣左右了。然後她指着我50美金買的鱷魚牌Polo衫，罵我說我「全身名牌卻上網寫網誌」，讓想要害她沒錢花。可想而知，她已經把我父母的錢當作她的錢了。而我不配花我父母的錢（實際上那件衣服還不到我月收入的千分之一）。

李麗麗對我釋放惡意，就比較微妙了。比如我高中二年級的時候，她和我班主任一起欺負我，逼我捐款幾千塊給學校。我一開始拒絕，她和班主任就煽動全班同學排擠我。（這件事也寫在這個網誌裏面。）

到了2021年的時候，我打開金山中學的學校宣傳，李麗麗還精準的把我的全國第二名金牌的名字，從金山中學光榮榜的名單上抹除掉。

在我身上，李麗麗一家把她們對我爸爸和媽媽的仇恨，釋放得淋漓盡致。也不用擔心影響她們從我父母的收入。

\ \\
李麗麗最仇恨的人，也是她的經濟來源。不得不從她最恨的人手裏不斷拿錢。這就是李麗麗糾結的人生吧。

\ \\
李新野\\
2025年7月19日

--------------------------------------------------------------------------

\

\textit{李新野網誌原文連結：https://sinyalee.com/blog/?p=1018}
